
\section{Graded rings and Hilbert functions}
\section{Systems of parameters and multiplicity}
\section{Dimension of extension rings}

\begin{theorem}[Theorem 13.2]\label{mcrt-auto-theorem-13-2}\dcref{13.2}\source{matsumura-commutative-ring-theory:page-0109}\uses{}
\begin{mathematicalrecord}
Let R = BnaoR, be a Noetherian graded ring with R,
Artinian, and let M be a finitely generated graded R-module. Suppose that
R= RO[xl,..., x,] with xi of degree di, and that P(M, t) is as above. Then
P(M, t) is a rational function of t, and can be written

        P(M, t) = f(t)/ fi (1 - tdi),
                        i=l

where f(t) is a polynomial with coefficients in Z.
\end{mathematicalrecord}
\end{theorem}
\begin{proof}
\begin{mathematicalrecord}
By induction on r?, the number of generators of R. When r = Q we
have R = R,, so that for n sufficiently large, M, = 0, and the power series
P(M, t) is a polynomial. When Y> 0, multiplication       by x, defines an Ro-
linear map M, - M, + d,; writing K, and L, +d,   for the kernel and cokernely
we get an exact sequence



Set K = OK,, and L = @ L,. Then K is a submodule of M, and L = MIxrM9
                                                                   tK
so that K and L are finite R-modules; moreover x,K = x,L = 0 SOtha

                                                                                 i
Graded rings, Hilbert function       and Samuel function                     95
913
and L can be viewed as R/x,R-modules,      and hence we can apply the
induction hypothesis to P(K, t) and P(L, t). Now from the above exact
sequence we get
        W,) - 4M?) + w, +d,) - 4L +d,) = 0.
If we multiply this by tnfdp and sum over n this gives
          tdrP(K, t) - tdrP(M, t) + P(A4, t) - P(L, t) = g(t),
where g(t)EZ[t].    The theorem follows at once from this. n
   A lot of information on the values of 1(M,) can be obtained from the
above theorem. Especially simple is the case d, = ... = d, = 1, so that R is
generated over R, by elements of degree 1. In this case P(M, t) = f(t)
(I - t)-?; if f(t) has (1 - t) as a factor we can cancel to get P in the form
          P(M,t)=f(t)(l        - t)-d with     ,f~Z[t],         d>O,
          and if d>O           then ,f(l)#O.
If this holds, we will write d = d(M). Since (1 - t)- 1 = 1 + t + t2 + . . . , we
can repeatedly differentiate both sides to get



(This can of course be proved in other ways, for example by induction on d.)
Iff(t)=a,+a,t+...+a,ts        then


(*I      W,) = a, (d~rr?)+a~(d:~~2)+...+a~(d+~~~-1);


                          = 0 for m < d - 1. The right-hand            side of (*) can be

formally rearranged as a polynomial        in n with rational coefficients, say cp(n);
then
                    f(l)
          ?(X)=(d- pXd-? I)!        + (terms of lower degree).

Since              coincides     with   the polynomial          m(m - 1). . . (m - d + 2)/
(d-   I)! for m >, 0, this implies the following      result.
\end{mathematicalrecord}
\end{proof}


\begin{theorem}[Theorem 13.3]\label{mcrt-auto-theorem-13-3}\dcref{13.3}\source{matsumura-commutative-ring-theory:page-0113}\uses{}
\begin{mathematicalrecord}
Let A be a semilocal Noetherian ring, and 0 + M? -
M -M?   + 0 an exact sequence of finite A-modules; then
         d(M) = max(d(M?),       d(M?))
If I is any ideal of definition of A, then XL - xfu? and xf\l, have the same
 leading coefficient.
Pro@. We can assume M? = MJM?. Then since M?/l*Mn = MJ(M? + I?M)
we have
         l(M/I?M)   = l(M/M? + I?M) + l(M? + I?M/I?M)
                    = I(M?/I?M?) + l(M?/M? n I?M).
Thus setting q(n) = I(M?/M?n        I?+ ? M), we have & = XL,, + 9. Since
moreover both &, and cptake on only positive values, d(M) coincides with
whichever is the greater of d(M?) and deg cp. However, by the Artin-Rees
lemma, there is a c > 0 such that
         n~~~l?~lM?cM?nl?~lMcI?~c~lM?,
and hence
         xL,(n) 2 q(n)>   x',,(n - c);
therefore cpand XL, have the same leading coefficient. n
    We now define a further measure 6(M) of the size of M: let 6(M) be the
smallestvalue of n such that there exist x1,. .. ,X,EIII for which l(M/x, M +
... + x,M) < a. When l(M) < rx, we interpret this as 6(M) = 0. If I is any
ideal of definition of A then l(M/IM) < co, so that 6(M) < number of
generators of I. Conversely, in the case that A is a local ring and M = A,
then 1(A/Z) < co implies that I is an m-primary ideal. Therefore in this case
6(A) is the minimum of the number of generators of m-primary ideals.
    We have now arrived at the fundamental theorem of dimension theory.
\end{mathematicalrecord}
\end{theorem}
\begin{proof}
\notready
\end{proof}


\begin{theorem}[Theorem 13.4]\label{mcrt-auto-theorem-13-4}\dcref{13.4}\source{matsumura-commutative-ring-theory:page-0113}\uses{mcrt-auto-theorem-6-4}
\begin{mathematicalrecord}
Let A be a semilocal Noetherian ring and M a finite
A-module; then we have
         dim M = d(M) = 6(M).
\end{mathematicalrecord}
\end{theorem}
\begin{proof}
\begin{mathematicalrecord}
Step 1. Each of d(M) and 6(M) are finite, but the finiteness of dimM
has not yet been established. First of all, let us prove that d(A) 3 dim A
for the caseM = A, by induction on d(A). Set m = rad (A). If d(A) = 0 then
l(A/m?) is constant for n >>0, so that for some n we have m? = m?+? and
by NAK, mn = 0. Hence any prime ideal of A is maximal and dim A = 0.
O+A/po    %4/p,       -B+O,

we have d(B) < d(A), and so by induction
         dimB<d(B)<d(A)-              1.
(The values of d(B) and dim B are independent of whether we consider B as
an A-module or as a B-module, as is clear from the definitions.) In B, the
image of p1 c ... c p, provides a chain of prime ideals of length e - 1, so that
         e- 1 ddimB<d(A)-          1;
hence e < d(A). Since this holds for any chain of prime ideals of A, this
proves dim A < d(A). For general M, by Theorem 6.4, there are submodules
MisuchthatO=M,cM,c...cM,=Mwith
         M,/M,-    1 E A/p,    and         p,eSpecA.
Since for an exact         sequence        O-+ M? -M      +      M?+O      of finite   A-
modules we have
         Supp(M)       = Supp(M?)uSupp(M?)

and
         dim M = max (dim M?, dim M?),
it is easy to see that
         d(M) = max {d(A/p,)}         3 max {dim(A/p,)}         = dim M.
 Step2. We show that 6(M) 3 d(M). IfG(M) = Othen l(M) < CCso that x,(n)
is bounded, hence d(M) = 0. Next suppose that 6(M) = s > 0, choose
Xl ,...,x,EITI such that l(M/x,M+.~~+x,M)<~,      and set Mi=M/X,M+
..* + XiM; then clearly 6(Mi) = 6(M) - i. On the other hand,
         ~(M,/v~?M,)      = I(M/x,M        + m?M)
                          = I(M/m?M)         - I(x,M/x,Mnm?M)
                          = I(M/m?M)        - I(M/(m?M:x,))
                          3 1(M/m?M)        - /(M/m?- ?M),
so that d(M,) > d(M) - 1. Repeating this, we get d(M,) 2 d(M) - s, but since
6(M,) = 0 we have d(M,) = 0, so that s 3 d(M).
   Step 3. We show that dim M 3 6(M), by induction on dim M. If dim M = 0
then Supp(M) c m-Spec A = b?(m) so that for large enough n we have
m? C ann M, and l(M) -C co, therefore 6(M) = 0. Next suppose that
dim M > 0, and let pi for 1 < i ,< t be the minimal prime divisors of ann (M)
with coht pi = dim M; then the pi are not maximal ideals, so do not contain
tn. Hence we can choose X~E~ not contained in any pi. Setting
M, = M/x,M we get dim M, < dim M. Therefore by the inductive hypo-
thesis 6(M,) 6 dim M,; but obviously 6(M) d @MI) + 1, so that 6(M) 6
dimM, + 1 <dimM.        n
\end{mathematicalrecord}
\end{proof}


\begin{theorem}[Theorem 13.5]\label{mcrt-auto-theorem-13-5}\dcref{13.5}\source{matsumura-commutative-ring-theory:page-0115}\uses{}
\begin{mathematicalrecord}
Let A be a Noetherian ring, and I = (a,, . . ,a,) an ideal
generated by Y elements; then if p is a minimal prime divisor of I we have
ht p < Y. Hence the height of a proper ideal of A is always finite.
\end{mathematicalrecord}
\end{theorem}
\begin{proof}
\begin{mathematicalrecord}
The ideal IA, c A, is a primary ideal belonging to the maximal
ideal, so that ht p = dim A, = 6(A,) < Y. w
\end{mathematicalrecord}
\end{proof}


\begin{theorem}[Theorem 13.6]\label{mcrt-auto-theorem-13-6}\dcref{13.6}\source{matsumura-commutative-ring-theory:page-0115}\uses{}
\begin{mathematicalrecord}
Let P be a prime ideal of height r in a Noetherian ring A.
Then
   (i) P is a minimal prime divisor of some ideal (a,, . . , a,) generated by
r elements;
   (ii) if b 1,. . . , b,EP we have ht P/(b,, . . . , b,) 3 r - s;
   (iii) if a,, . . . , a, are as in (i) we have
            htP/(a,,...,u,)=r-i            for ldibr.
\end{mathematicalrecord}
\end{theorem}
\begin{proof}
\begin{mathematicalrecord}
(i) A, is an r-dimensional local ring, so that by Theorem 4 we can
choose r elements a,, . . . , u,EPA~ such that (a,, . , u,)A, is PAP-primary.
Each ui is of the form an element of P times a unit of A,, so that without
loss of generality we can assume that a+P. Then P is a minimal Prime
divisor of (a,, . . . , u,)A.
   (ii) Set A = A/(b,, . . . , b,), B = P/(bI , . . , b,) and ht P = t. Then by (i),
                 Graded rings, Hilbertfunction     and Samuel function                   101
     813
     there   exist   Cl,..., C,EP    such that P is a minimal prime divisor of
     (El,..., t ;)A. Then    P  is a  minimal prime divisor of (b,, . . . , b,, cl,. . . ,cJ,
     and hence Y d s + t by Theorem 5.
        (iii) The ideal P/(al , . . , a i) is a minimal prime divisor of (a,, 1,. . , &)
     kA/(al,...,     ai), hence ht P/(al, . , ai) d r - i. The opposite inequality was
     proved in (ii). W
\end{mathematicalrecord}
\end{proof}


\begin{theorem}[Theorem 13.7]\label{mcrt-auto-theorem-13-7}\dcref{13.7}\source{matsumura-commutative-ring-theory:page-0116}\uses{mcrt-auto-theorem-5-6}
\begin{mathematicalrecord}
Let A = @nBOAn be a Noetherian graded ring.
         (i) If I is a homogeneous ideal and P is a prime divisor of I then P is also
     homogeneous.
         (ii) If P is a homogeneous prime ideal of height Y then there exists a
     sequence P = P, 1 P, 3 ... 3 P, of length r consisting of homogeneous
     prime      ideals.
     J+YK$ (i) P can be expressed in the form P = ann (x) for a suitable element
     x of the graded A-module A/Z. Let UEP, and let x = x0 +x1 + ... + x,
     anda=a,+a,+,            + ... + a, be decompositions into homogeneous terms.
     Then since ax = 0,
                apxO = 0,   apxl +a,+,x,=O,          a,x,+a,+,x,       +ap+2x0=0,...,
     from which we get six, = 0, six, = 0,. . . , and finally a;+?~ = 0. It follows
      that aif EP, but since P is prime, apEP. Thus ap+l + ... + a,EP, so
    ?that in turn ap+ l~P. Proceeding in the same way, we see that all the homo-
     geneous     terms of a are in P, so that P is a homogeneous ideal.
        (ii) First of all note that we can assume that A is an integral domain.
     To see this, if we take a chain P = p0 2 ... 1 p, of prime ideals of length
     r then p, is a minimal prime divisor of (0), and so by (i) is a homogeneous
     ideal; so we can replace A by A/p,. Now choose a homogeneous element
     0 # b, EP; then by Theorem 6, ht (P/b, A) = r - 1, and so there is a minimal
    .prime divisor Q of b, A such that ht (P/Q) = r - 1; since Q # (0), it is a
     height 1 homogeneous prime ideal. By the inductive hypothesis on r
     applied to P/Q there exists a chain P = P, I P, I ... I> P,- 1 = Q of homo-
  ? geneousprime ideals of length r - 1, and adding on (0) we get a chain
    ,Of length r. w
,./
 :.     Let us investigate more closely the relation between local rings and
;- @aded rings.
-!l Theorem 13.8. Let k be a field, and R = k[<,,. . ., ~$1 a graded ring
1%. generated by elementstl,. . ,<, of degree 1; set M = c 5iR, A = R, and
*.**.
pgj mEMA.
?%     (i) Let x be the Samuel function of the local ring A, and cpthe Hilbert
   .function of the graded ring R; then q(n) = x(n) - x(n - 1);
  j (ii)dimR=htM=dimA=degcp+l;
      (iii) gr,(A) N R as graded rings.
\end{mathematicalrecord}
\end{theorem}
\begin{proof}
\begin{mathematicalrecord}
M is a maximal        ideal of R so that
            m?/m?+?     N Mn/Mni     1 2: R,;
hence x(n) - x(n - 1) = l(mn/m?+l) = 1(R,) = q(n), and so dim A = deg x = 1 +
deg cp. Then since A = R,, we have dim A = ht M. After this it is enough
to prove that dim R = ht M. First of all, assume that R is an integral
domain, so that by Example 3 in the section on Hilbert functions and by
Theorem 5.6, we have
          l+degcp>tr.deg,R=dimR>,htM;
putting this together with ht M = dim A = 1 + deg cp, we get dim R = ht M.
Next for general R, let P,, . . . , P, be the minimal prime ideals of R; then by
Theorem 7, these are all homogeneous ideals, and each R/Pi is a graded
ring. Choosing P, such that dim R = dim R/P, and using the above result,
we get
         dim R = dim R/P, = ht M/P, < ht M ,< dim R,
so that dim R = htM as required. We have R, c M? c m? with m?/m?+ ? 1
R,, and so taking an element x of R, into its image in m?/m?+ ? we obtain
a canonical one-to-one map R 1gr,,,A,    and it is clear from the definition
that this is a ring isomorphism.  n
\end{mathematicalrecord}
\end{proof}


\begin{theorem}[Theorem 13.9]\label{mcrt-auto-theorem-13-9}\dcref{13.9}\source{matsumura-commutative-ring-theory:page-0117}\uses{mcrt-auto-theorem-15-7}
\begin{mathematicalrecord}
Let (A, m, k) be a Noetherian local ring, and set G = gr,,A;
then dim A = dim G.
\end{mathematicalrecord}
\end{theorem}
\begin{proof}
\begin{mathematicalrecord}
Letting q be the Hilbert polynomial         of G, we have dim A =
 1 + deg cp (by Theorem 4), and by the previous theorem this is equal to
dim G. n
    In fact the following more general theorem holds: for I a proper ideal in a
Noetherian local ring A, set G = gr,(A); then dim A = dim G. This will be
proved a little later (Theorem 15.7).
\end{mathematicalrecord}
\end{proof}


\begin{theorem}[Theorem 13.10]\label{mcrt-auto-theorem-13-10}\dcref{13.10}\source{matsumura-commutative-ring-theory:page-0119}\uses{mcrt-auto-theorem-13-4}
\begin{mathematicalrecord}
Eagon). Let A be a Noetherian ring and M be an r x s
matrix (I d s) of elements of A. Let Ir be the ideal of A generated by the t x t
minors of M. If P is a minimal prime divisor of I, then we have
           ht P d (r - t + l)(s - t + 1).
\end{mathematicalrecord}
\end{theorem}
\begin{proof}
\begin{mathematicalrecord}
Induction on r. When Y = 1 we have t = 1, and so (r - t + l)(s - t +
 1) = s. The ideal I, is generated by s elements, so that the assertion is just
the principal ideal theorem (Theorem 5) in this case. Next assume that r > 1.
Localising at P we may assume that A is a local ring with maximal ideal P,
and that I, is P-primary.
    If t = 1, then I, is generated by rs elements and (r - t + 1) (s - t + 1) = rs,
so our assertion holds also for this case. Therefore we assume t > 1. If at
least one of the elements of M is a unit of A, then by what we said above, I, is
generated by (t - 1) x (t - 1) minors of a (r - 1) x (s - 1) matrix, and again
we are done. Therefore we assume that all the elements of M are in P. Now
comes the brilliant idea. Let M? be the matrix with elements in B = A[X]
obtained from M by replacing a, 1 by a,, + X, and let I? be the ideal of B
generated by the t x t minors of M?. Since t > 1 and aijgP for all i and j we
have I? c PB. We also have I? + XB = Z,B + XB since both sides have the
same image in B/XB = A. Therefore PB is a minimal prime divisor of I? by
the lemma. Since the element a, 1 + X of M? is not in PB, we have
ht PB < (r - t + l)(s - t + 1) by our previous argument. Since ht PB = ht P,
as we can see by Theorems 4 and 5, we are done. n

      14   Systems of parameters    and multiplicity
            Let (A, m) be an r-dimensional Noetherian local ring; by Theorem
13.4, there exists an m-primary ideal generated by r elements, but none
generated by fewer. If a,,. . . , a,~nt generate an m-primary             ideal,
Iu  I,?.., a}  is  said to  be   a system    of parameters  of  A  (sometimes
abbreviated to s.0.p.). If M is a finite A-module with dim M = s, there exist
Yl,..., y,Em such that l(M/(y, ,..., y,)M) < cc, and then (y, ,..., y} is
said to be a system of parameters of M.
   If we set A/m = k, the smallest number of elements needed to generate
m itself is equal to rank,m/m?;       (here rank, is the rank of a free module
over k, that is the dimension of m/m2 as k-vector space). This number is
called the embedding dimension of A, and is written emb dim A. In general
           dim A < emb dim A,
\end{mathematicalrecord}
\end{proof}


\begin{theorem}[Theorem 14.1]\label{mcrt-auto-theorem-14-1}\dcref{14.1}\source{matsumura-commutative-ring-theory:page-0120}\uses{mcrt-auto-theorem-13-6}
\begin{mathematicalrecord}
Let (A, m) be a Noetherian local ring, and x1,. . . , x, a system
  of parameters. Then
      (i) dim A/(x,, . . . , xi) = r - i for 1 < i < r.
      (ii) although it is not true that ht (x1,. . . , Xi) = i for all i for an arbitrary
  system of parameters, there exists a choice of x1,. . . , x, such that every
  subset Fc(xr,...,          x,] generates an ideal of A of height equal to the
 number of elements of F.
\end{mathematicalrecord}
\end{theorem}
\begin{proof}
\begin{mathematicalrecord}
(i) is contained in Theorem 13.6. We now prove the second half of
 (iii. If r< 1 the assertion is obvious; suppose that r > 1. Let poj (for 1 <
l<ee) be the prime ideals of A of height 0. Choosing xrom not contained
 in any poj, we have ht (xi) = 1. Next letting plj (for 1 <j 6 e,) be the minimal
 prime divisors of (xi), SO that ht Plj = 1, and choosing xZEm not contained
 in any poj or any plj, we have ht(xJ = 1, ht(x,,x,) = 2; if r = 2 we?re
 done. If r > 2 we choose x3Em not contained in any minimal prime
 divisor of (0), (x,), (x,), (x1, x,), and proceed in the same way to obtain
 the result.
      We now give an example where ht (x1,. . , xi) < i. Let k be a field and
 set R = k[X, Y, Zj; let I = (X)n( Y, Z), and write A = R/Z, and x, y, z for
 the images in A of X, Y, Z. The minimal prime ideals of A are (x) and
 TV,z); now A/(X) = R/(X) N k[lx Z] is two-dimensional                    and A/(y, z) N
 R/( y, Z) N k[X]      .is one-dimensional,    so that dim A = 2. {y, x + Z} is a
 system of parameters of A; in fact xy = xz = 0. so that x2 =x(x + a)~
(.Y,x + Z) and z* = z(x + z)~(y,x + z). However, y is contained in the
.minimal prime ideal (y, z) of A, and hence ht (y) = 0. n
\end{mathematicalrecord}
\end{proof}


\begin{theorem}[Theorem 14.2]\label{mcrt-auto-theorem-14-2}\dcref{14.2}\source{matsumura-commutative-ring-theory:page-0120}\uses{}
\begin{mathematicalrecord}
Let (R, m) be an n-dimensional                    regular local ring, and
 Xl , . . . , Xi elements of m. Then the following conditions are equivalent:
     (1) Xl,... ,Xi is a subset of a regular system of parameters of R;
     (2) the images in m/m* of x r , . . . , xi are linearly independent over R/m;
     (3) R/(x,,... ,xJ is an (n - i)-dimensional regular local ring.
 ~00f.(1)~(2)1fx,,...,Xi,Xi+1,...           , x, is a regular system of parameters then
 their images generate m/m* over k = R/m, and since rank,m/m* = n they
 must be linearly independent over k.
     (I)*(3)              We know that dim R/(x l,...,xi)=n-ii,      and the images of
x1+1, . . . , x, generate the maximal ideal of R/(x,, . . . , xi).
     (3)=>(l) If the maximal ideal m/(x1,. . . , xi) of R/(x,, . . . , xi) is generated
bY the images of y I,..., y,-iEm                then m is generated by x1,. ..,xi,
:Y1,...,yn-i.
\end{mathematicalrecord}
\end{theorem}
\begin{proof}
\notready
\end{proof}


\begin{theorem}[Theorem 14.3]\label{mcrt-auto-theorem-14-3}\dcref{14.3}\source{matsumura-commutative-ring-theory:page-0121}\uses{mcrt-auto-theorem-11-2}
\begin{mathematicalrecord}
A regular local ring is an integral domain.
\end{mathematicalrecord}
\end{theorem}
\begin{proof}
\begin{mathematicalrecord}
Let (R,m) be an n-dimensional regular local ring; we proceed by
induction on n. If n = 0 then m is an ideal generated by 0 elements, so
that m = (0). This in turn means that R is a field. Thus a zero-dimensional
regular local ring is just a field by another name.
   When n = 1, the maximal ideal m = xR is principal and ht m = 1, so that
there exists a prime ideal p # m with m 3 p. If yen we can write y = xa with
aER, and since x$p we have aEp; hence p = xp, and by NAK, p = (0). This
proves that R is an integral domain. (There is a slightly different proof in the
course of the proof of Theorem 11.2; as proved there, a one-dimensional
regular local ring is just a DVR by another name.)
   When n> 1, let pi,... ,pI be the minimal prime ideals of R; then since
m $ m2 and m + pi for all i, there exists an element xEm not contained in
anyofm2,p,,.     . . , p,. (see Ex. 1.6). Then the image of x in m/m? is non-zero,
so that by the previous theorem R/xR is an (n - 1)-dimensional regular
local ring. By the induction hypothesis, R/xR is an integral domain, in
other words, xR is a prime ideal of R. If pi is one of the minimal prime ideals
contained in xR then since x$pi, the same argument as in the n = 1 case
shows that p1 = xp,, and hence p1 = (0). n
\end{mathematicalrecord}
\end{proof}


\begin{theorem}[Theorem 14.4]\label{mcrt-auto-theorem-14-4}\dcref{14.4}\source{matsumura-commutative-ring-theory:page-0121}\uses{}
\begin{mathematicalrecord}
Let (A, m, k) be a d-dimensional       regular local ring; then
          gr,,V) 2: kCX,, . . .,XJ,
and if x(n) is the Samuel function of A then
                   n+d
          x(n)=      d          for all n 3 0.
                     (     >
\end{mathematicalrecord}
\end{theorem}
\begin{proof}
\begin{mathematicalrecord}
Since m is generated by d elements, gr,(A) is of the form
4X i , . . . , X,1/1, where I is a homogeneous ideal. Now if I # (0) let SEI be a
non-zero homogeneous element of degree r; then for n > r the homogeneous
piece of k[X]/Z       of degree n has length         at most

                     which is a polynomial       of degree d - 2 in n. This
implies that the Samuel function of A is of degree at most d - 1, and
contradicts dim A = d. Hence I = (0); the second assertion follows from
the first. n
   Let (A, m) be a Noetherian local ring. Elements yl,. . . ,y,Em are said to be
analytically independent if they have the following property; for every
\end{mathematicalrecord}
\end{proof}


\begin{theorem}[Theorem 14.5]\label{mcrt-auto-theorem-14-5}\dcref{14.5}\source{matsumura-commutative-ring-theory:page-0122}\uses{mcrt-auto-theorem-13-4}
\begin{mathematicalrecord}
Let (A,m) be a d-dimensional Noetherian local ring and
  xr,.i.,x,, a system of parameters of A; then x1,. . .,xd are analytically
  independent.
\end{mathematicalrecord}
\end{theorem}
\begin{proof}
\begin{mathematicalrecord}
Set q = xxiA. Since q is an ideal of definition of A, by
   Theorem 13.4, xi(n) = l(A/q?) is a polynomial of degree d in n for n >>0. Set
  ,@t = k; we say that a homogeneous form f (X)Ek[X1,.      . . ,X,] of degree n
  is a null-form    of q if F(x,....x,)~q?m     for any homogeneous form
  F(X)EACXl, . . ., X,] which reduces to f(X) modulo m. Write n for the ideal
  of&X,,...,   X,] generated by the null-forms of q. Then
              k[xlln = 0 q?/q?m,
  and writing cp for the Hilbert polynomial            of k[X]/n, we have q(n)
  = I(q?/q?m) for n >>0. The right-hand side is just the number ofelements in a
  minimal basis of q?, so that cp(n).l(A/q) >, l(q?/q?+?).   Now
              4q?lq??)         = x34   - x?dn - 1)
  is a polynomial in n of degree d - 1, so that deg q > d - 1, but if n # (0) this
  is impossible. Thus n = (0), and the statement in the theorem follows at
  once. n

  lirurtiplicity
   &et (A, m) be a d-dimensional Noetherian local ring, M a finite A-module,
   and q an ideal of definition of A (that is, an m-primary ideal). As we saw
   ti $13, the Samuel function ,!(M/q?+? M) = x$(n) can be expressed for
:. 5~0 as a polynomial in n with rational coefficients, and degree equal to
   dim M, and therefore at most d. In addition, this polynomial can only take
   tnteger values for n >>0, so it is easy to see by induction on d (using the fact
   that x(n + 1) - x(n) has the same property) that

              Xg,(n) = grid + (terms of lower order),

b.%th eeZ. This integer will be written e(q, M). By definition       we have the
'3dlOwing          property.
 :>A.=
C?
.::.Fom~la 14.1. e(q, M) = n+m
                             lim cl(M/q?M),   and in particular, if d = 0 then
                                  n*
$jyk M) = Z(M).
       rom this we see easily the following:
      rmula 14.2. e(q, M) > 0 if dim M = d, and e(q, M) = 0 if dim M < d;
     rmula 14.3. e(q?, M) = e(q, M)rd;
Formula 14.4. If q and q? are both m-primary                   ideals and q 3 q? then
eh M) d 44, Ml.
   We set e(q,A) = e(q),and define this to be the multiplicity of q. In addition,
we will refer to the multiplicity e(m) of the maximal ideal as the multiplicity
of the local ring A, and sometimeswrite e(A) for it. For example, if A is a
regular local ring then by Theorem 4, we can see that e(A) = 1.
\end{mathematicalrecord}
\end{proof}


\begin{theorem}[Theorem 14.6]\label{mcrt-auto-theorem-14-6}\dcref{14.6}\source{matsumura-commutative-ring-theory:page-0123}\uses{}
\begin{mathematicalrecord}
Let O+ M? -M                -    M?+O       be an exact sequence of
finite A-modules. Then
        e(q, Ml = e(q,M?) + e(cr,M?).
\end{mathematicalrecord}
\end{theorem}
\begin{proof}
\begin{mathematicalrecord}
We view M? as a submodule of M. Then
         /(M/q? M) = I(M?/q? M?) + I(M?/M? n q? M),
and obviously q?M? c M?nq?M. On the other hand by ArtinRees, there
exists c > 0 such that
         M?nq?Mcq       n-CM?        for all   y1> C.
Hence
         /(Ml/q?-?M?)   < l(M?/M?n       q?M) < l(M?/q?M?).
From this and Formula 14.1 it follows easily that

         e(q, M) - e(q,M?) = lim f$(M?/M?n              q?M)   = e(q, M?).
                             n-sn
\end{mathematicalrecord}
\end{proof}


\begin{theorem}[Theorem 14.7]\label{mcrt-auto-theorem-14-7}\dcref{14.7}\source{matsumura-commutative-ring-theory:page-0123}\uses{}
\begin{mathematicalrecord}
Let (pl,. . . ,p} be all the minimal prime ideals of A such
that dim A/p = d; then



where si denotesthe image of q in A/p, and 1(M,) standsfor the length of M, as
A,-module.
Proqf (taken from Nagata [Nl]). We write CJ= &l(M,) and proceed by
induction on o. If CJ= 0 then dim M < d, so that the left-hand sideis 0, and
the right-hand side is obviously 0; now supposec > 0. Now there is some
p~{pi,. . ,p} for which M, #O; then p is a minimal element of
Supp(M). Hence p~Ass(M), that is M contains a submodule N isomorphic
to A/p. Then
        e(q, W = 4% N) + e(s, MIN.
On the other hand, N, N Ap/pA, and N,: = 0 for pi # p, so that 1(N,) = 1,
and the value of g for M/N has decreasedby one, so that the theorem holds
for M/N. However, from the definition
         4% N) = e(q, A/p) = e(9,A/p), where 9 = (q + P)/P.
Putting this together, we seethat the theorem also holds for M. l
  Theorem 7 allows usto reduce the study of e(q, M) to the casethat A isan
T&orem 14.9. Let (A, m) be a Noetherian local ring, q an ideal of definition
   ofA,andx,,...,    xd a system of parameters of A contained in q. Suppose that
   xiEqv* for 1 < i 6 d. Then for a finite A-module M and s = 1,. . . ,d we have
             e(q/h ,..., x,),M/(xl ,..., x,)M)3v,v,...v,e(q,M).
   In particular if s = d, we have
               I(M/(x,,      . ,qJM) 3 v1v2.. . vde(q,M).
\end{mathematicalrecord}
\end{theorem}
\begin{proof}
\begin{mathematicalrecord}
It is enough to prove the case s = 1. We set A? = A/x, A,q? = q/x, A,
   &f=M/x,M            and v=vl.      By Theorem 1, we have dimA?=d-              1. On the
   other hand,
               l(M?/q??M?)= l(M/x, M + q?M)
                          = l(M/q?M) - 1(x, M + q?M/q?M).
     In addition, in view of (x1 M + q?M)/q?M 2: x1 M/x, Mnq?M                        N Ml
     (q?M:x,) and q?-?M c qnM:xl, we have
  .?           - 1(x, M + q?M/q?M) 3 - l(M/q?-?M),
 ?I and   therefore
     I
               &Vf'/q'"M')     > l(M/q?M)    - l(M/q?-?M).
 i;, When n >>0 the right-hand side is of the form
-7-t
   ?/       eh Ml
 d<
;p.?        T[nd        - (n - v)~] + (polynomial of degree d - 2 in n)

                   = pv.nd-l           + (polynomial     of degree d - 2 in n),

        that the assertion is clear. n
       A case of the above theorem which is particularly          simple, but important,


         eorem 14.10. Let (A, m) be a d-dimensional Noetherian local ring, let
             xd be a system of parameters of A, and set q = (x1,. . . ,x,); then

       d if in addition       xiEm? for all i then l(A/q) 2 v?e(m).

           orem 14.11. Let A, m, xi and q be as above. Let M be a finite A-module,
           set A? = A/x, A, M? = M/x, M and q? = q/x1 A = c$xiA?. Then ifx, is a
non-zero-divisor      of M, we have the following                   equality
          4-r, M) = 44, W.
Proof.   Since l(M?/q??+ 1M?) = l(M/x, M + q?+ 1 M) we have
          1(M/q? + r M) - l(M?/q?n+ r M?)=I(x,M+q?+?M/q?+?M)
            = l(x,M/x,Mnq?+?M)         = l(M/(q?+?M:x,))
            = l(M/q?M) - l((q?+?M:x,)/q?M).
On the other hand, setting a = x$xiA we have q = x,A + a and q?+? =
xlqn+a?+l,      and therefore
           qntl M:x, = q?M + (a?+?M:x,).
Moreover, by Artin-Rees, there is a c > 0 such that for n > c we have
a?+lMnx,M=a?-c(ac+?        Mnx,M),     and therefore a?+rM:xl c an-CM.
Thus
         (q?+?M:x,)/q?M   = (q?M +(a?+?M:x,))/q?M
                          c (q?M + a?-?M)/q?M
                          ~a?-?M/a?-?Mnq?M.
NOW an-?M/a?-?M n q?M is a module over A/q?, and since a is generated
                                            n-c+d-2
by d - I elements, an-c is generated by                   elements. Thus
                                                d-2     >
for n > c we have



where m is the number of generators of M. The right-hand side is a
polynomial of degree d - 2 in II, so that
          eh?, M ?) = (d - l)! lim l(M?/qln+ ? M?)/n?- ?
                               n-cc
                    = (d - l)! lim [l(M/q?+?M)                 - l(M/q?M)]/nd-?
                               n-30
                    = e(q, M). H
\end{mathematicalrecord}
\end{proof}


\begin{theorem}[Theorem 14.12]\label{mcrt-auto-theorem-14-12}\dcref{14.12}\source{matsumura-commutative-ring-theory:page-0125}\uses{}
\begin{mathematicalrecord}
Lech?s lemma). Let A be a d-dimensional Noetherian
local ring, and x r , . . , ,xd a system of parameters of A; set q = (x1,. . .,xd), and
suppose that M is a finite A-module. Then
                                    l(M/(x;?,         . , x,y?)M)
          e(q, M) =           lim
                      min(vi)+m             VI...Vd
\end{mathematicalrecord}
\end{theorem}
\begin{proof}
\begin{mathematicalrecord}
If d = 0 then both sidesare equal to l(M). Ifd = 1 then the right-hand
side is exactly Formula 14.1 which defines e(q, M). For d > 1 we use
induction on d.
   Setting Nj= {meMlxjm=O)          we have N, c N, c . .. so that there is a
c > 0 such that N, = N,, 1 = .*.. If we set M? = xt M then x, is a non-zero-
divisor for M?, and there is an exact sequence0 + N, -M           -M?     40.
Since N, is a module over A/x?, A we have dim N, < d, and therefore
e,q, &I) = e(q, M?). On the other hand,
            l(M/(x;?, . . . ,xdYd)M) - l(M?/(x;?,.    . . ,xF)M?)
                = l(N, + (xy, . . ,xp)M/(x;?,           . ,xi?)M)
                = l(N,/N, n (x;? , . . , x;?)M)
                d l(N,/(x;?, . . . ,x;?)N,).
  If v1 7 c then x;l N, = 0, and N, is a module over the (d - 1)-dimensional
  Local ring A/x;A, so that by induction there is a constant C such that as
  min (vi) + CC we have
             l(N,l(x;?, . >x2) NJ = l(N,/(x?,Z, . . . , x;?) N,) < CT,. . . vd.
  Therefore,
            lim [l(M/(xT?, . . . , xi?)M) - l(M?/(xr?, . . . , x~?)h?f?)]/v,. . . vd = 0.
  This means that we can replace M by M? in the theorem, and so we can
  assume that x1 is a non-zero-divisor in M. Then by the previous theorem we
  have e(q, M) = e(q, J@),with q = q/x, A and ~ = M/x, M. If we furthermore
  set
           E=(xT,...,   xF)M      and F = M/E
  then by Theorem 9, we have
           e(q,   M).v, . ..vd < l(M/(x?,?,       . ,xi?)M)     = l(F/x;?F)

                               = i$ 1(x;- l F/x?,F)           ~v,~(F/x,F)=v,~(M/x,M+E)
                               = VI l(~/(X~)       . . ,xJy)Al).
   Then by induction on d we have
          lim l(M/(x;l,. . . , x;?)M)/v,.
                                       . . vd = lim l(R/(x;?, . . . , xdyd)fl)/vz.. . vd
                                              = e(q, M).     n

     Although we will not use it in this book, we state here without proof a
  remarkable result of Serre which shows that multiplicity can be expressed
  as the Euler characteristic of the homology groups of the Koszul complex
  (discussed in 5 16).
\end{mathematicalrecord}
\end{proof}


\begin{theorem}[Theorem 14.13]\label{mcrt-auto-theorem-14-13}\dcref{14.13}\source{matsumura-commutative-ring-theory:page-0127}\uses{}
\begin{mathematicalrecord}
Let (A, m) be a Noetherian local ring, q an m-primary ideal
and b a reduction of q; then b is also m-primary, and for any finite A-module
M we have
         e(q, W = 46 Ml.
\end{mathematicalrecord}
\end{theorem}
\begin{proof}
\begin{mathematicalrecord}
If q?+l = bq? then q?+? c b c q, hence b is also m-primary.
Moreover,
         l(M/b?+?M)         3 I(M/q?+?M)   = l(M/b?q?)   3 f(M/b?M),
so that e(q, M) = e(b, M) follows easily.
\end{mathematicalrecord}
\end{proof}


\begin{theorem}[Theorem 14.14]\label{mcrt-auto-theorem-14-14}\dcref{14.14}\source{matsumura-commutative-ring-theory:page-0127}\uses{}
\begin{mathematicalrecord}
Let (A, m) be a d-dimensional Noetherian local ring, and
suppose that A/m is an infinite field; let q = (u,, . . . , u,) be an nt-primary
ideal. Then if yi = xaijuj for 1 < i < d are d ?sufficiently general? linear
combinations of u Ir.. . , us, the ideal b = (y 1,. . . , yd) is a reduction of q and
 {yl,. . , yd} is a system of parameters of A.
\end{mathematicalrecord}
\end{theorem}
\begin{proof}
\begin{mathematicalrecord}
If d = 0 then q? = (0) for some Y > 0, hence (0) is a reduction of q
so that the result holds. We suppose below that d > 0.
    Step 1. Set A/m = k and consider the polynomial ring k[X,, . . . , X,] (or
k[X] for short). For a homogeneous form q(X) = I&X,, . . . ,X,)EA[X]               of
degree n, we write @(X)Ek[X]         for the polynomial obtained by reducing
the coefficients of q modulo m. As in the proof of Theorem 5 we say that
@(X)Ek[X]       is a null-form of q if cp(u1,. . . ,u,)Eq?m; this notion depends
not just on q, but also on ul,. .., u,. However, for fixed @ it does not
depend on the choice of cp.We write Q for the ideal of k[X] generated by
all the null-forms of q, and call Q the ideal of null-firms of q. One sees
easily that all the homogeneous elements of Q are null-forms of q, and
that the graded ring k[X]/Q has graded component of degree n isomorphic               )
to qn/qnm, so that we have
         KXIIQ = @q"/q"m = gr,(40A,,k.
                      n,O

Write q(n) for the Hilbert function of k[X]/Q; then
          44 = 4q?lq?m) f h?/q?+ ?1 G cpW Wd
(see the proof of Theorem 5). We know that for n >>0, the function l(q?/q??)
is a polynomial in n of degree d - 1 (where d = dim&            Thus from the
above inequality, cp is also a polynomial          of degree d - 1, so that by
Theorem 13.8, (ii), we have dimk[X]/Q         = d.
   Now set I/ = ct kX,, and let P,, . . . , p, be the minimal prime divisors Of
i=l

Hence we can take a linear form 1,(X)~l/ not belonging to any Pi. If
d 71 then similarly we can take ~,(X)EV such that 12(X) is not contained
m any minimal prime divisor of (Q, II(X)), and, proceeding in the same
way, we get II(X), . . .,&(X)E V such that (Q, 1,, . . . ,1,) is a primary ideal
belonging to (XI, . . . ,X,).
   Step 2. We let b be the ideal of A generated by linear combinations
                                                                   t


L,(U) = C       (for 1 < i < of u I , . . , u, with coefficients in A. Then if we
              Uijuj                     t)


f@t li(X)=ti(Xj    =CaijXj3   a necessary and sufficient condition for b to
be a- reduction of qis -that the ideal (Q, 1,). . . ,1,) of k[X] is (X,, . . . ,X,)-
primary.
Proof of necessity. Suppose that bqr = q?+l. Then if M = M(X ) is a
monomial of degree r + 1 in X,, . . . ,X,, we can write



where the F,(X) are homogeneous                    forms of degrees r with coefficients in
A. Thus
        ii?(X) - c li(X)Fi(X)gQ.
Hence

         (X l,...,Xs)rfl~(Q,ll,...,lt).
Proof of sufficiency.                We go through     the same argument    in reverse: if
~~--~I,F,EQ    then
         M(U)-              CL~(U)F~(U)E~?+~ITI,
so that qr+l c bqr + q?+?nt; thus by NAK, q*+? = bq?.
    Step 3. Putting together Steps 1 and 2 we see that q has a reduction
b=Cyl,... ,y,J generated by d elements. Both q and its reduction b are
m-primary ideals, so that y,, . . . ,y, is a system of parameters of A. We
are going to prove that there exists a finite number of polynomials D,(Zij)
for 1< c(6 v in sd indeterminates Z, (for 1 < i d d and 1 <j < s) such that
d linear combinations yi = C aijuj (for 1 < i < d) generate a reduction ideal
of 9 if and only if at least one of D,(aij) # 0. (The expression ?d sufficiently
general linear combinations? in the statement of the theorem is quite vague,
hut in the present case it has a precise interpretation as above,)
   Let G,(X) , . . . ,G,(X) be a set of generators of Q, with Gj homogeneous
of degree e> For any sd elements aij of k (for 1 < i f d and 1 <j < s), set
hw) = c aijxi. w e write I, for the homogeneous component of degree n
of a homogeneous ideal I c k[X,, . . . ,X,1, and in particular               we write
(X 1,. . . ,X,), = V,, so that
            w 1I..., XJc(Q,4   ,..., Id)0Vn=(Q,11,...,Id)n.
Set c, = dim, V,. We have

          (Q, 1l,...,Z~,={CliFi+        CGjHjlFiEV,-,           and    HjEV/,-,j}.

Let K1,..., K, be the elements obtained as xliFi + c GjHj as the Fi run
through a basis of V,-, and the Hj run independently through a basis
of Vnpe,; it is clear that they span (Q, I,, . . . ,I& Each of K,, . . . ,K, is a
linear combination of the c, monomials of degree n in the Xi, with linear
functions in the rxij as coefficients; we write out these coefficients in a
c, x w matrix. If qp,,,(aij) for 1 d v < p, are the c, x c, minors of this matrix
then the necessary and sufficient condition for (X,, . . . , X,)n c (Q, I,, . . . , lJ to
hold is that at least one of the 4pn(aij) is non-zero. Therefore the ideal
(Q,h,..., Id) will fail to be (X,, . . . , X,)-primary if and only if the
quantities uij satisfy (~?,,(a~~)= 0 for all n and all v. However, the ring k[Zij]
is Noetherian, so that the ideal of k[Zij] generated by all of the cp,,(Zij) is
generated by finitely many elements D,(Zij) for 1 f CId u. These D, clearly
meet our requirements.         n
\end{mathematicalrecord}
\end{proof}


\begin{theorem}[Theorem 15.1]\label{mcrt-auto-theorem-15-1}\dcref{15.1}\source{matsumura-commutative-ring-theory:page-0131}\uses{}
\begin{mathematicalrecord}
Let cp:A -B             be a homomorphism     of Noetherian rings,
 and P a prime ideal of B; then setting p = Pn A, we have
   (i) ht P d ht p + dim B&B,;
   (ii) if q is flat, or more generally if the going-down theorem holds between
A and B, then equality holds in (i).
\end{mathematicalrecord}
\end{theorem}
\begin{proof}
\begin{mathematicalrecord}
We can replace A and B by A, and B,, and assume that (A, m) and
(B, n) are local rings, with mB c n. Rewriting (i) in the form
            dim B < dim A + dim B/mB
makes clear the geometrical content. To prove this, take a system of
parameters x1,. . . , x, of A, and choose yr,. . . , y,eB such that their images
in B/mB form a system of parameters of B/mB. Then for v, p large enough
 we have ny c mB + z y,B and rn@c xxjA, giving ny? c x y,B + C XjB.
Hence dim B 6 r + s.
   (ii) Let dim B/mB = s, and let n = P, 3 P, I ... 3 P, be a strictly decreas-
ing chain of prime ideals of B between n and mB. Obviously we have
P,nA=mforO~i~s.NowsetdimA=randletm=p,~p,~~~~~p,
be a strictly decreasing chain of prime ideals of A; by the going-down
theorem, we can construct a strictly decreasing chain of prime ideals of B
            Ps3Ps+lx?-,3Ps+r            such that P,, i n A = pi.
Thus dim B 2 r + s, and putting this together with (i) gives equality.         n
\end{mathematicalrecord}
\end{proof}


\begin{theorem}[Theorem 15.2]\label{mcrt-auto-theorem-15-2}\dcref{15.2}\source{matsumura-commutative-ring-theory:page-0131}\uses{}
\begin{mathematicalrecord}
Let q:A -+ B be a homomorphism        of Noetherian rings?
and suppose that the going-up theorem holds between A and B. Then if
p and q are prime ideals of A such that p 3 q, we have
        dim B @ K(P) 2 dim B 0 K(q).
\end{mathematicalrecord}
\end{theorem}
\begin{proof}
\begin{mathematicalrecord}
Set r = dimB@Ic(q)     and s = ht(p/q). We choose a strictly
\end{mathematicalrecord}
\end{proof}


\begin{theorem}[Theorem 15.3]\label{mcrt-auto-theorem-15-3}\dcref{15.3}\source{matsumura-commutative-ring-theory:page-0132}\uses{}
\begin{mathematicalrecord}
Let cp:A+B be a homomorphism    of Noetherian rings, and
      suppose that the going-down theorem holds between A and B. If p and q
      are prime ideals of A with p 1 q then
              dim B@c(p) <dim B@K(q).
\end{mathematicalrecord}
\end{theorem}
\begin{proof}
\begin{mathematicalrecord}
We may assume that ht(p/q)= 1, and it is enough to prove that,
      given a chain P, cP, c . . . c P, of prime ideals of B lying over p such that
      ht(P,/P, _ r) = 1 we can construct a chain of prime ideals Q. c Q 1 c . . . c Q,
      of B lying over ~7such that
              QicPi    (Odidr)   and ht(Qi/Qi-l)=l       (O<i<r).
      We can find Q. by going down. If r > 1 then take x~p -q and let ??, , . . . ,T,
      be the minimal prime divisors of Q,+xB.        Then ht(TJQ,)= 1, while
      ht(PJQ,) > 2, hence we can choose


      Let Qr be a minimal prime divisor of Qo+ yB contained in P,. Then
      ?ht(Qr/Q,J = 1, and Q1 # q for all i, hence c$x)#Q,.
         Therefore Q,nA#p,    and since ht(p/q)= 1 we must have QlnA=q.    By
      the same method we can successively construct Q1 , Qz, . . . , Q,. q




      2. polynomial and formal power seriesrings
\end{mathematicalrecord}
\end{proof}


\begin{theorem}[Theorem 15.4]\label{mcrt-auto-theorem-15-4}\dcref{15.4}\source{matsumura-commutative-ring-theory:page-0132}\uses{mcrt-auto-theorem-31-7}
\begin{mathematicalrecord}
Let A be a Noetherian        ring, and X,, . . . , X, indeterminates
     over A. Then
              dimA[X,,...      ,X,]=dimA[X,         ,..., X,j=dimA+n.
  + %of. It is enough to consider the case n = 1. For any p&pecA,                      the
:I?. hjt A[x]    &K(P)    = K(p)[x]    is a principal ideal ring, and therefore one-
2?:
dimensional; also A[X] is a free A-module, hence faithfully flat, so by
Theorem 1, (ii), dim A[X] = dim A + 1.
   For ,4[XJ it is not true in general that A[Xa@,rc(p)    and tc(p)[X]
coincide; however, if m is a maximal ideal of A we have
         AUxljOlc(m)=ABx40(A/m)=(A/m)[Cxn,
and this fibre ring is one-dimensional.    Also, as we saw on p. 4, every
maximal ideal ?%I of A[XJ is of the form W = (m,X), where m = sJJln~
is a maximal ideal of A. Thus for a maximal ideal 9.R of ,4%X] we have
         ht?9JI = ht(mnnA) + 1;
conversely, if m is a maximal ideal of A then ht(m, X) = ht m + 1, and
putting these together gives dim A [Xl = dim A + 1. m
Remark 2. It is not necessarily true that a maximal     ideal of A[X] lies
over a maximal ideal of A. For example, if A is a DVR and t a
uniformising element then A[t-?1 = K is the field of fractions of A, so
that A[X]/(tX - 1) 2: K, and (tX - 1) is a maximal ideal of A[X]; however,
(tX- l)nA=(O).
Remark 2. It is quite common         for fibre rings of A -+A[X,,.    . .,X,1
to have dimension strictly greater than n. For example, let k be a field and
set A = k[Y,Z]. It is well-known that the field of fractions of k[Xj has
infinite transcendence degree over k(X) (see[ZS], vol. II, p. 220). Let u(X),
u(X)~k[Xj      be two elements algebraically independent over k(X), and
define a k-homomorphism      (continuous for the X-adic topology)


bycp(X!=X,cp(Y)=u(X),cp(Z)=u(X).IfwesetKercp=PthenPnA=(O),
and A[XJ/P z k[XJ is one-dimensional.         Now every maximal ideal of
A[Xj has height 3, and, as we will see later, A[Xj is catenary, so that
ht P=2. Thus we see that the generic libre of A-+A[Xj            is two-
dimensional.

3. The dimension inequality
We say that a ring A is universally catenary if A is Noetherian and every
finitely generated A-algebra is catenary. Since any A-algebra generated
by n elements is a quotient of A[X,, . , . , X.1, and since a quotient of a
catenary ring is again catenary, a necessary and sufficient condition for a
Noetherian ring A to be universally catenary is that A[X,, . . . ,X,] is
catenary for every n 2 0. (In fact it is known that it is sufficient for A[XI]
to be catenary, compare Theorem 31.7.)
\end{mathematicalrecord}
\end{theorem}
\begin{proof}
\notready
\end{proof}


\begin{theorem}[Theorem 15.5]\label{mcrt-auto-theorem-15-5}\dcref{15.5}\source{matsumura-commutative-ring-theory:page-0133}\uses{}
\begin{mathematicalrecord}
I. S. Cohen [3]). Let A be a Noetherian   integral domain, and
(*) htP + tr.deg,(,, K(P) d htp + tr.deg, B,
   where tr.deg,B is the transcendence degree of the field of quotients of B
   over that of A.
\end{mathematicalrecord}
\end{theorem}
\begin{proof}
\begin{mathematicalrecord}
We may assume that B is finitely generated over A. For if the right
   hand side is finite and m and t are non-negative integers such that m < htP
   and t < tr.deg,(,+(P),       then there is a prime ideal chain
   p=po3       P, X3??X P, in B. Take ai~Pi-Pi+l,          06i<m,      and let
   C1,...,~,~B    be such that their images modulo P are algebraically
   independent over A/p. Set C = A[{a}, (cjj]. If the theorem holds for C,
   then we have m + t < htp + tr.deg, C < htp + tr.deg, B. Letting m and t vary
   we seethe validity of (*).
      We may furthermore assume,by induction, that B is generated over A by
   a single element: B = A[x]. We can replace A by A, and B by B, = AJx],
   and hence assumethat A is local and p its maximal ideal. Set k = A/p and
   w&e B = A[X]/Q. If Q = (0) then B = A[X] and by Theorem 1 we have
   htP = htp + ht(P/pB), and since B/pB = k[X] we have either P = pB or
   ht(P/pB) = 1. In both casesthe equality holds in (*).
      If Q # (0) then tr.deg, B = 0. Since A is a subring of B we have
   &A=(O),        so that writing K for the field of fractions of A we have
   htQ = htQK[X] = I. Let P* be the inverse image of P in A[X]. Then
   P=P*/Q,K(P)=K(P*),           and     htP<htP*-htQ=htP*-l=htp+l-
   tr.deg,&P*) - 1= htp -tr.deg,(,,lc(P).     n
\end{mathematicalrecord}
\end{proof}


\begin{theorem}[Theorem 15.7]\label{mcrt-auto-theorem-15-7}\dcref{15.7}\source{matsumura-commutative-ring-theory:page-0137}\uses{}
\begin{mathematicalrecord}
Let A be a Noetherian ring and I a proper ideal; then
setting R = R(A, I) and G = gr,(A) we have
         dimR=dimA+         1, dimG<dimA.
If in addition A is local, then
         dim G = dim A.
\end{mathematicalrecord}
\end{theorem}
\begin{proof}
\notready
\end{proof}
